\section{核心原则：从候选线索到校准证据}

本文采用三层机制表述：候选线索、可靠证据、图像条件化参照。这个表述把局部响应、证据选择和测试时校准放在各自的位置上。

\paragraph{候选线索。}
CLIP 初始异常概率 $s_i^0$、局部纹理变化、边界突变和高频响应都属于候选线索。它们提示某个区域与异常判定相关，同时也是混合观测量：缺陷、正常纹理和物体结构都会触发它们。因此，候选线索回答“哪里值得重新检查”，可靠证据回答“哪里相对于当前图像参照异常”。

\paragraph{可靠证据。}
可靠证据要求候选线索通过当前图像上下文的约束。对正常证据而言，一个位置应具有较低的初始异常响应，并呈现稳定的局部结构；对异常证据而言，一个位置应具有较高的初始异常响应，并呈现结构边界之外的局部变化。换言之，证据选择需要同时考虑语义倾向和局部可靠性。

\paragraph{图像条件化参照。}
参照由测试图像内部的可靠证据形成评分基准。该参照可以表现为校准后的 normal/abnormal prototypes，也可以表现为对局部评分函数的图像条件化约束。本文采用前者：通过可靠 patch evidence 构造 visual prototypes，并用它们对 CLIP text prototypes 做保守校准。

这一参照概念与已有路线形成清晰分工。WinCLIP+ 的参照来自少量外部 normal reference images；VCP-CLIP 的图像条件化来自经过跨数据集训练的 Pre-/Post-VCP 模块；TPT 和 CoCoOp 说明视觉语言模型中测试时或图像条件化 prompt adaptation 已经存在。本文聚焦更窄的问题：在目标类别参考库缺失、prompt 和 adapter 参数保持固定的设定下，单张测试图内部的哪些局部响应可以作为 prototype calibration evidence。

这个原则导出三个可证伪预期。第一，直接把局部可靠性图当作异常图会把正常纹理和结构边界纳入同一响应，因此它构成负控。第二，基于 CLIP 初始语义响应选择校准证据会继承初始响应的偏差；语义之外的局部可靠性约束应当改善证据选择的稳定性。第三，保守校准应当控制异常侧错误证据对 prototypes 的污染；正常图误报面积提供相应检验。第~\ref{sec:mechanism-tests} 节的消融表和图正是围绕这些预期组织。
